% ===========================================================
% Conference Paper Starter Template
% ===========================================================
% Replace the content below with your own title, authors,
% abstract, sections, and citations.
% ===========================================================

\documentclass[11pt]{article}

% ---- Packages ----
\usepackage[T1]{fontenc}
\usepackage[utf8]{inputenc}
\usepackage[a4paper,margin=1in]{geometry}
\usepackage{natbib}
\usepackage{amsmath,amssymb,amsfonts}
\usepackage{graphicx}
\usepackage{textcomp}
\usepackage{xcolor}
\usepackage{hyperref}
\usepackage{setspace}

\graphicspath{{figures/}}
\onehalfspacing

% ===========================================================
\begin{document}

\title{The Corporate Stage: Algorithmic Surveillance, Hybrid Visibility, \\
and the Paradox of Spectatorship in the Modern Workplace}

\author{
  Vivien Jiaqian Zhu \\
  \textit{Department of East Asian Languages and Cultures, University of California, Berkeley} \\
  \textit{Visiting Scholar, Stanford University} \\
  \textit{Research Fellow, University of Cambridge} \\
  \texttt{email@example.com}
}

\date{}

\maketitle

% ===========================================================
\begin{abstract}
% TODO: 150--250 word abstract for IntechOpen submission.
% Should state: (1) the problem -- the convergence of algorithmic
% monitoring and performative self-presentation in contemporary
% hybrid/remote workplaces; (2) the theoretical apparatus -- Goffman's
% dramaturgical model of everyday performance, Foucault's account of
% disciplinary and panoptic power, and Zuboff's theory of surveillance
% capitalism; (3) the central claim -- that hybrid visibility produces
% a "paradox of spectatorship," in which workers are simultaneously
% audience and performer, watched and self-watching, under conditions
% that exceed the classical panopticon; (4) what the chapter
% contributes to workplace studies / organizational theory.
\end{abstract}

\textbf{Keywords:} algorithmic surveillance, hybrid work, dramaturgical
theory, panopticism, surveillance capitalism, spectatorship, workplace
visibility

% ===========================================================
\section{Introduction}
\label{sec:intro}

% NOTE: opening paragraphs drafted below; revise once empirical
% examples / statistics are sourced and verified.

A worker logs into a video call from a spare bedroom, angles the
camera to exclude an unmade bed, and opens a second window running
software that logs keystrokes, application switches, and idle time.
Somewhere on a dashboard she will likely never see in full, these
signals resolve into a productivity score. She is, in this moment,
watched by a colleague on the call, watched by software she cannot
see operating, and watching herself -- adjusting posture, tone, and
visible activity in anticipation of both audiences at once. This
scene, unremarkable in its ordinariness across contemporary hybrid
and remote workplaces, is the empirical starting point for this
chapter's argument: that the shift to hybrid and algorithmically
monitored work has not simply added new tools to an old practice of
supervision, but has reorganized the basic conditions under which
workers are seen, and see themselves, at work.

Workplace visibility has never been neutral. What has changed is its
architecture. Where earlier regimes of supervision relied on a
supervisor's line of sight, a factory floor's layout, or a manager's
sporadic presence, contemporary hybrid work distributes visibility
across software that operates continuously, asymmetrically, and
largely outside the worker's control. The question this chapter
takes up is not whether workers are watched -- they always have
been -- but how the specific combination of algorithmic monitoring
and hybrid work reorganizes the relationship between watching and
being watched, producing a condition this chapter calls the
\textit{paradox of spectatorship}.

\subsection{The Corporate Stage as Analytic Frame}

This chapter borrows, and deliberately strains, Erving Goffman's
dramaturgical vocabulary. Goffman's front stage and back stage
described a world in which performances were bounded in space and
time: a worker performed professionalism on the office floor and
relaxed once out of view. Hybrid work collapses this architecture.
The home becomes a front stage that a webcam or activity monitor can
activate at any moment, and the back stage -- the space in which a
worker could once be off duty, off script, unobserved -- shrinks
correspondingly or disappears. The \textit{corporate stage}, as this
chapter develops the term, names a performance space that is no
longer bounded by walls or shifts but is instead switched on and off
by software, with the worker frequently unable to know which
condition currently obtains. This is a context Goffman did not, and
could not, theorize: one in which the audience is not only human but
algorithmic, not only present but recording, and not only observing
but scoring.

\subsection{Contributions / Chapter Roadmap}
The chapter makes the following moves:
\begin{itemize}
  \item Synthesizes Goffman, Foucault, and Zuboff into a single
        framework for reading hybrid workplace visibility.
  \item Names and develops the ``paradox of spectatorship'' as a
        distinct condition of contemporary algorithmic management.
  \item Considers what this reorganization of visibility implies for
        worker autonomy and wellbeing, and for how organizations
        might design monitoring practices that do not collapse the
        distinction between observation and total exposure.
\end{itemize}

% ===========================================================
\section{Theoretical Framework}
\label{sec:framework}

\subsection{Goffman: The Presentation of Self as Front-Stage Labor}

Goffman's dramaturgical model treats social life as a series of
performances staged for particular audiences, sustained through what
he called impression management: the continuous, often unconscious
work of controlling how one appears to others~\citep{goffman1959}.
Central to this model is the distinction between front stage, where
a performance is delivered, and back stage, where a performer can
drop the performance, rehearse, or simply rest. Crucially, Goffman
located the back stage in physical space -- a break room, a car, a
home -- precisely because performances required somewhere to end.
Hybrid work erodes this spatial guarantee. When the home itself
becomes a monitored workspace, the back stage does not merely shrink;
it becomes conditional, switched on and off by whether a camera is
active or a monitoring application is running. Workers report
adapting home environments -- posture, lighting, even the visible
contents of a room -- not for a single scheduled performance but for
a front stage that could, in principle, be activated at any moment.
Impression management, in this setting, becomes a standing
condition of labor rather than an episodic task bounded by shift or
meeting.

\subsection{Foucault: Panopticism and the Internalization of the Gaze}

Foucault's account of the panopticon turns on a specific mechanism:
disciplinary power that operates not through constant observation
but through the \emph{uncertainty} of being observed, such that the
watched party internalizes the gaze and disciplines their own
conduct regardless of whether anyone is actually watching at a given
moment~\citep{foucault1975}. This uncertainty principle is what
makes the panopticon efficient -- surveillance does not need to be
continuous to be effective, only possible. Algorithmic workplace
monitoring both intensifies and departs from this model. It
intensifies it insofar as software can, in principle, observe
continuously, removing the gaps in which a watched party could
once assume privacy by default. But it also departs from Foucault's
model in a specific way: where the panopticon relies on
\emph{uncertainty} to produce self-discipline, many monitoring
systems make their presence explicit -- a recording indicator, a
dashboard the worker can also see -- converting internalized
discipline into an openly acknowledged, sometimes actively managed,
condition. The worker under algorithmic monitoring is often not
guessing whether they are watched, but calibrating behavior against
a known, quantified record of having been watched.

\subsection{Zuboff: Surveillance Capitalism and Behavioral Data Extraction}

Zuboff's theory of surveillance capitalism describes an economic
logic in which human experience is unilaterally claimed as raw
material, translated into behavioral data, and rendered into
predictive products sold in markets the person generating the data
never sees or consents to in any meaningful
sense~\citep{zuboff2019}. Zuboff developed this account primarily
through consumer platforms -- search, social media, smart devices --
but the underlying logic extends with little modification into the
workplace. Keystroke logs, activity timestamps, meeting attendance
patterns, and productivity scores constitute a comparable behavioral
surplus, extracted not from consumers but from labor itself, and
often justified internally as a byproduct of legitimate
productivity or security tools rather than named as data extraction.
The asymmetry Zuboff identifies as central to surveillance
capitalism -- one party accumulates knowledge and predictive power,
the other supplies the raw behavioral material largely without
recourse -- maps closely onto the relationship between an employee
and the monitoring infrastructure of a hybrid workplace.

\subsection{Toward a Synthesis: Hybrid Visibility}

Read together, these three accounts describe three distinct but
compounding operations: Goffman locates the site of performance and
shows how that site has lost its guaranteed off-switch; Foucault
explains the disciplinary mechanism by which observation, actual or
possible, shapes conduct; and Zuboff supplies the economic logic
that gives monitoring its institutional purpose -- behavioral data
is not merely observed but captured, aggregated, and rendered
actionable. \textit{Hybrid visibility}, as this chapter uses the
term, names the condition produced when these three operations occur
together: a performance space without a reliable back stage
(Goffman), sustained by infrastructure that watches continuously
rather than merely possibly (an intensification of Foucault), for
the purpose of producing extractable behavioral data (Zuboff).
% TODO: connect to relevant empirical/educational-technology
% literature on dashboards, learning analytics, and visibility as a
% design problem -- e.g.,~\citep{rets_TODO}. Citation still needs
% verification; do not cite in submitted draft until confirmed.
None of the three theorists anticipated this specific convergence,
which is precisely what makes their combination useful: each
supplies a piece the others lack.

% ===========================================================
\section{Hybrid Visibility in the Modern Workplace}
\label{sec:hybrid}

\subsection{From Panopticon to Platform: What Changes}
% TODO: Distinguish classical disciplinary visibility (architectural,
% spatial, uncertain) from platform-based visibility (continuous,
% quantified, data-retentive, algorithmically scored).

\subsection{The Home as Front Stage}
% TODO: The domestic space reconfigured as a site of professional
% performance; blurred back-stage/front-stage boundaries under
% webcam-based and activity-tracking monitoring regimes.

\subsection{Metrics, Scores, and the Datafied Self}
% TODO: Productivity dashboards, "active time," algorithmic scoring
% as a new register of visibility distinct from direct human
% observation.

% ===========================================================
\section{The Paradox of Spectatorship}
\label{sec:paradox}

\subsection{Audience and Performer at Once}
% TODO: Core theoretical move of the chapter -- the worker as
% simultaneously spectator of their own datafied self (via dashboards,
% self-monitoring apps) and performer for an algorithmic/managerial
% audience. Unpack the reflexive loop this creates.

\subsection{Illustrative Cases}
% TODO: 2--3 concrete illustrative cases or vignettes (e.g., call
% center monitoring, remote-work productivity software, gig-platform
% rating systems) to ground the theoretical argument. Confirm whether
% IntechOpen chapter format expects a dedicated "case studies" section
% or integrated examples.

\subsection{Limits of the Paradox: Resistance and Opacity}
% TODO: Worker tactics of partial opacity, gaming of metrics, or
% refusal -- as a counterpoint to a purely totalizing surveillance
% narrative.

% ===========================================================
\section{Discussion}
\label{sec:discussion}

% TODO: Implications for organizational design, management practice,
% and worker wellbeing. Position the chapter's contribution within the
% edited volume's broader theme (\textit{The Evolving Landscape of
% Workplace -- Challenges and Opportunities}, ed. C\u{a}t\u{a}lina Radu).

% ===========================================================
\section{Conclusion}
\label{sec:conclusion}

% TODO: Restate the paradox of spectatorship as the chapter's central
% contribution; gesture toward future research (e.g., comparative
% regulatory approaches, worker-side sousveillance tools).

% ===========================================================
\bibliographystyle{apalike}
\bibliography{references}

\end{document}
